\documentclass[12pt]{article}
\usepackage[margin=0.75in]{geometry}
\usepackage{fontspec}
\setmainfont{Latin Modern Roman}
\usepackage{microtype}
\usepackage{fancyhdr}
\usepackage{titlesec}
\usepackage{enumitem}
\usepackage{xcolor}
\usepackage[hidelinks]{hyperref}
\setlength{\parindent}{1.1em}
\setlength{\parskip}{0.35em}
\setlength{\headheight}{15pt}
\titleformat{\section}{\large\bfseries\scshape}{}{0pt}{}
\titleformat{\subsection}{\normalsize\bfseries}{}{0pt}{}
\pagestyle{fancy}
\fancyhf{}
\fancyhead[L]{Whately Elements of Logic}
\fancyhead[R]{Teacher Guide}
\fancyfoot[C]{\thepage}
\renewcommand{\headrulewidth}{0.4pt}

\begin{document}
\begin{center}
{\LARGE\bfseries Teacher Guide: Week 4}\par\vspace{0.25em}
{\large\scshape Because: How Syllogisms Work}\par\vspace{0.4em}
{12th Grade Long Logic Unit Based on Richard Whately's \textit{Elements of Logic}}\par\vspace{0.6em}
\rule{0.82\textwidth}{0.4pt}
\end{center}

\section*{Week 4 Overview}
Week 4 moves from clear propositions to practical syllogism structure. Students learn that a syllogism uses two premises to support a conclusion through a shared middle term. The week intentionally stays practical: students identify the conclusion, major term, minor term, middle term, major premise, and minor premise without yet memorizing traditional mood-and-figure classifications. The literary anchor uses \textit{Julius Caesar} to show how a serious argument can become dangerous when a debatable middle term makes fear sound like proof.

\section*{Essential Question}
When do two claims actually force a conclusion?

\section*{Student-Facing Materials}
\begin{itemize}[leftmargin=*]
\item \textit{Because: How Syllogisms Work} — student reading handout.
\item \textit{Week 4 Argument Lab: Practical Syllogism Maps} — student lab for mapping premises, conclusion, major/minor/middle terms, and producing a portfolio window note.
\item \textit{Julius Caesar Lit-Anchor: The Serpent's Egg} — Shakespeare anchor passage for testing Brutus's preemptive argument against Caesar and Cassius's persuasive comparison argument.
\end{itemize}

\section*{Learning Goals}
Students will be able to:
\begin{enumerate}[leftmargin=*]
\item Define a syllogism as a deductive argument using two premises to support a conclusion.
\item Identify the conclusion of a practical syllogism.
\item Identify the major term and minor term from the conclusion.
\item Identify the middle term that appears in the premises but not in the conclusion.
\item Distinguish major premise from minor premise.
\item Explain why the middle term functions as the bridge or medium of proof.
\item Recognize when a middle term fails to prove the conclusion.
\item Map a hidden or implied syllogism from ordinary speech or literature.
\item Apply syllogism mapping to Brutus's and Cassius's arguments in \textit{Julius Caesar}.
\end{enumerate}

\section*{Key Vocabulary}
syllogism, premise, conclusion, term, major term, minor term, middle term, major premise, minor premise, deduction, valid form, medium of proof

\section*{Core Teaching Emphasis}
The main danger this week is that students will treat a syllogism map as a truth machine. Emphasize the distinction between \textit{following from premises} and \textit{having true premises}. A syllogism map can expose the structure of Brutus's argument, but it does not automatically prove Caesar is dangerous. Students should learn to say: ``If these premises are true and the terms connect, the conclusion follows; but we still have to test the premises.''

\section*{Weekly Lesson Sequence}

\subsection*{Day 1: From Propositions to Arguments}
\begin{itemize}[leftmargin=*]
\item Review Week 3: propositions have subject, predicate, quality, and quantity.
\item Introduce the shift from one proposition to two propositions supporting a conclusion.
\item Work through the Socrates example from the reading.
\item Students identify premise and conclusion in several simple classroom examples.
\end{itemize}

\subsection*{Day 2: Major, Minor, and Middle Terms}
\begin{itemize}[leftmargin=*]
\item Teach the three terms from the conclusion outward: minor term as subject of the conclusion, major term as predicate of the conclusion, middle term as the shared bridge in the premises.
\item Use the ``unjust acts'' example from the reading.
\item Complete Part I and the first item of Part II in the Argument Lab as a class.
\item Emphasize that the middle term must actually connect the other terms, not merely appear in both premises.
\end{itemize}

\subsection*{Day 3: When the Bridge Fails}
\begin{itemize}[leftmargin=*]
\item Use the ``All careful thinkers ask questions'' example to show a tempting but invalid movement.
\item Have students complete the remaining Part II examples in pairs.
\item Ask students to explain failures in ordinary language before using technical labels.
\item Close with the question: what did the argument prove, and what did it only seem to prove?
\end{itemize}

\subsection*{Day 4: Julius Caesar Lit-Anchor}
\begin{itemize}[leftmargin=*]
\item Read Brutus's Act 2, Scene 1 ``serpent's egg'' soliloquy from the lit-anchor.
\item Build the guided syllogism map on the board: major premise, minor premise, conclusion, major term, minor term, middle term.
\item Focus discussion on the weak point: Brutus has not proved Caesar is already ruled by passion; he argues from what Caesar may become.
\item Read Cassius's Act 1, Scene 2 passage as the independent comparison.
\item Students begin the \textit{Julius Caesar Practical Syllogism Map} portfolio artifact.
\end{itemize}

\subsection*{Day 5: Portfolio Map and Synthesis}
\begin{itemize}[leftmargin=*]
\item Students finish either the Argument Lab window note or the lit-anchor practical syllogism map.
\item Hold a short seminar: Which argument is more logically disciplined, Brutus's or Cassius's? Which premise is most debatable?
\item Connect back to Whately: a middle term is the medium of proof, but only if it truly connects the premises to the conclusion.
\item Exit writing: ``A syllogism map helps a reader because...'' Require one example from Whately, Brutus, or Cassius.
\end{itemize}

\section*{Answer Guidance for Brutus}
One strong map of Brutus's reasoning:
\begin{itemize}[leftmargin=*]
\item \textbf{Major premise:} Anyone who will probably become a dangerous tyrant if given supreme power should be stopped before he gains that power.
\item \textbf{Minor premise:} Caesar will probably become a dangerous tyrant if crowned.
\item \textbf{Conclusion:} Caesar should be stopped before he is crowned.
\item \textbf{Major term:} people who should be stopped before gaining power.
\item \textbf{Minor term:} Caesar.
\item \textbf{Middle term:} one who will probably become a dangerous tyrant if given supreme power.
\end{itemize}

The minor premise is the main pressure point. Brutus admits that he has not personally known Caesar's passions to overpower his reason. His argument relies on a general pattern about ambition rather than direct proof about Caesar. Students should see the difference between a warning, a probability, and a proven conclusion.

\section*{Answer Guidance for Cassius}
Cassius's argument is less disciplined but rhetorically powerful. One possible map:
\begin{itemize}[leftmargin=*]
\item \textbf{Major premise:} No man who is merely equal to Rome's nobles should tower over Rome like a king.
\item \textbf{Minor premise:} Caesar is merely equal to Brutus and the other nobles.
\item \textbf{Conclusion:} Caesar should not tower over Rome like a king.
\item \textbf{Middle term:} merely equal to Rome's nobles.
\end{itemize}

This map exposes a gap. Cassius may show that Caesar is physically weak or not naturally superior to Brutus, but that does not by itself prove Caesar is politically dangerous, tyrannical, or deserving of death. Cassius's speech is useful because it mixes argument with resentment, honor, comparison, and inherited Roman identity.

\section*{Common Student Errors}
\begin{itemize}[leftmargin=*]
\item Confusing the middle term with whichever word appears most often.
\item Identifying the conclusion before checking the exact subject and predicate.
\item Treating vivid imagery, such as the serpent's egg, as proof.
\item Saying an argument is invalid when they actually mean one premise is doubtful.
\item Treating Brutus as simply right or wrong before mapping his reasoning.
\end{itemize}

\section*{Connection to Future Weeks}
Week 5 can now introduce mood, figure, and formal validity because students have practiced the parts of a syllogism in ordinary language. Week 6 will then distinguish validity from truth and soundness. The Brutus case should be saved as a recurring example: his argument may be formally organized, but the truth of its premises remains contested.

\end{document}
